\documentclass[11pt]{article}
\usepackage{amsmath,amssymb}
\usepackage[margin=2.5cm]{geometry}
\title{Learning Across Domains}
\author{Ada Lovelace}
\date{}

\begin{document}
\maketitle

\begin{abstract}
We study how a classifier trained on one distribution behaves on another, and give a simple bound on the target error.
\end{abstract}

\section{Introduction}
Machine learning models are often deployed in conditions that differ from training. We recieve data from a source domain $\mathcal{S}$ but are evaluated on a target domain $\mathcal{T}$.

\section{Main result}\label{sec:main}
\paragraph{The theory.} The $\mathcal H\Delta\mathcal H$-divergence measures how well a pair of hypotheses can tell the two domains apart.

\begin{equation}\label{eq:bound}
  \epsilon_{\mathcal T}(h) \le \epsilon_{\mathcal S}(h) + \tfrac12\, d_{\mathcal H\Delta\mathcal H}(\mathcal S,\mathcal T) + \lambda .
\end{equation}

Equation~\eqref{eq:bound} shows that a small divergence implies good transfer.

\begin{itemize}
  \item \textbf{Bound:} source error plus divergence.
  \item \emph{Intuition:} similar domains transfer well.
\end{itemize}

\section{Conclusion}
Section~\ref{sec:main} gives the bound; tighter versions are left for future work.
\end{document}
